% =====================================================================
\section{Introduction}
\label{sec:intro}

Dense visual correspondence (optical flow, geometric and semantic matching)
relies heavily on synthetic training data, because dense ground truth is scarce
on real images
\cite{mayer2016FlyingThingsDataset,sintel,greff2021kubric,pointodyssey}. Most
progress targets \emph{how} to render this data (fidelity, lighting, speed)
\cite{greff2021kubric,pointodyssey}, on the assumption that closing the
appearance gap drives transfer. Yet non-photorealistic data often transfers well
\cite{Baradad2021LookingAtNoise,white2009mandelbulb}, and a logically prior
question gets far less attention: given a controllable pipeline, \emph{what}
should it render for a given target? As rendering becomes more efficient and the
appearance gap narrows, what to generate becomes the binding constraint.

We find that a synthetic dataset's value is not intrinsic but depends on the
target: the best source for one target is middling for another, so ``best
dataset'' is meaningful only relative to deployment. Matching a source's
distribution to a target's is also inherently \emph{directional}. In a shared
motion-descriptor space a source can place mass on motions the target never uses
(\textbf{off-target mass}; precision-analogous) or fail to cover the motions it
needs (\textbf{coverage}, or missing support; recall-analogous); we measure these
as the two directed nearest-neighbor distances between the datasets' flow-vector
distributions. Symmetric discrepancies such as FID
\cite{heusel2018ganstrainedtimescaleupdate} and sliced Wasserstein average the
two directions, whereas generative-model evaluation separates them for sample
quality \cite{sajjadi2018assessing_precision_recall,kynkaanniemi2019improved_pr}.
For \emph{transfer}, we find the separation matters: most of the predictive
signal sits in one direction.

\begin{quote}
\textbf{The coverage principle.} \emph{The most reliable single predictor of
transfer is \emph{coverage}: how well the source's motion covers the target's,
measured by the directed distance $\dbt$. Off-target mass, motion the source renders but the target never uses,
contributes comparatively little, so it is usually better to \emph{cover} the
target's motion than to \emph{match} its full distribution.}
Two consequences recur: symmetric and distribution distances (FID, Wasserstein)
trail coverage because they fold the less-useful direction back in; and the
appearance analogue of coverage carries little target-selective signal, setting a
source's overall quality rather than which target it suits. Throughout,
``pretrained'' refers to the \emph{image-feature backbone} (\eg, ImageNet
ResNet-101), not correspondence-task pre-training; the correspondence head is
always trained on the synthetic source.
\end{quote}

We study this on a transfer study of 142 trained correspondence models: four
architecture families (CATs++ \cite{cho2022cats++}, RAFT \cite{teed2020raft},
GLU-Net \cite{truong2020glunet}, and FlowFormer \cite{huang2022flowformer}) in 13
backbone configurations over 11 training sources, evaluated on 9 benchmarks
($1{,}278$ transfer measurements). We probe it from several angles: replication
across three distance estimators; a negative control in appearance-feature space;
a construct-validity check, in which matching the descriptor alone recovers each
target's expected motion regime; and a controllable grid crossing motion against
appearance, on which a target-tuned source beats a generic baseline. A controlled
motion-magnitude grid maps the coverage response (an inverted U, in which both
under-coverage and over-shoot cost transfer); a cleaner test that isolates
coverage at fixed off-target mass remains future work.

We make the following contributions:
\begin{enumerate}
\item \textbf{A directional coverage signal for correspondence transfer}
  (Sec.~\ref{sec:law}): of the distances we test, motion coverage ($\dbt$) is the
  most reliable single predictor of transfer (within-context $\rho{\approx}0.57$,
  about $92\%$ of the reproducibility ceiling), holding across architectures with
  each reported in its designed training regime (Sec.~\ref{sec:setup}).
\item \textbf{A render-free selection rule, zero- and few-shot}
  (Secs.~\ref{sec:predict}--\ref{sec:absolute}): rank candidate sources by motion
  coverage $\dbt$, with no target-side run, no labels, no learned weights, and no
  differentiable renderer. Zero-shot it picks a source within a median $0.6$
  \pck{} of the best (versus $15.1$ at random, median regret) and orders
  clearly-separated candidates correctly $79\%$ of the time, close to an
  independent cross-architecture retraining ($82\%$). Absolute \pck{} is set
  largely by a source-level realism effect the motion distances cannot see; given
  a handful of observed runs as anchors, the same score recovers it few-shot (MAE
  ${\approx}10$ with the benchmark observed, widening to ${\approx}30$ when both
  source and benchmark are held out) under leakage-safe splits. The anchors fix the absolute scale; coverage fixes the order.
\item \textbf{Cover, do not match: off-target mass is nearly free}
  (Secs.~\ref{sec:law}--\ref{sec:predict}): off-target mass adds little predictive
  signal ($\dtb$, $\rho{\approx}0.03$), so the symmetric and distribution
  distances (FID, sliced Wasserstein-2) trail coverage by folding it back in, and
  the appearance analogue of coverage carries little target-selective signal. The
  practical reading is to \emph{add} coverage of the target's motion rather than
  \emph{match} its full distribution, which helps explain why broad generic sets
  such as FlyingThings transfer well, and how a tuned source can improve on them.
\item \textbf{Interventional support} (Sec.~\ref{sec:interv}): a render-free
  generator search, scoring each candidate on the flow it produces with no
  photorealistic rendering and no training in the loop and scored only by motion
  coverage, recovers each target's expected motion regime (a forward camera dolly
  for KITTI, object motion for FlyingThings). A source tuned this way matches or
  beats a generic synthetic baseline in all twelve \emph{KITTI} settings (two
  benchmarks $\times$ three architectures $\times$ two backbones), and on a
  controllable grid crossing motion against appearance, motion rather than
  appearance is the target-specific carrier of transfer.
\end{enumerate}

The stratification and design partition were fixed before FlowFormer was added
to the study; its four configurations entered as a held-out fourth architecture
family and fell in line with the coverage signal on real-motion targets without
adjustment (pretrained $\dbt$ $+0.70$/$+0.67$, Table~\ref{tab:law}); on semantic
targets it is a weaker but still positive coverage signal, where a symmetric
distance is competitive (Sec.~\ref{sec:law}). Either way it entered the analysis
unchanged, a check that the analysis choices are not fit to the data they explain.
